\section{Experimental Setup}
\label{sec:setup}

\subsection{Datasets}

\begin{table}[!t]
\caption{Datasets and evaluation scope (all public, Zenodo-pinned).}
\label{tab:datasets}
\centering
\small
\begin{tabular}{@{}p{0.30\linewidth}p{0.30\linewidth}p{0.26\linewidth}@{}}
\toprule
Property & RAN PM counters~\cite{lehoczky2026ran} & CTTC slicing~\cite{farreras2024cttc} \\
\midrule
Role & Temporal demand + Stages 1, 4, 5 & Slice-level policy comparison (Stages 2--3) \\
Raw size & 302{,}052 rows, 75 sectors & 250 simulation snapshots \\
Granularity & 15-min slots & Steady-state snapshots \\
Evaluation series & 1{,}778 contiguous slots ($\approx$18.5 days) & 229 retained snapshots \\
Capacity $C$ & 143{,}954.5 (1.3$\times$ observed peak) & Per-sample link capacity \\
Split & 70/15/15 chronological & 70/30 by sample id \\
Test scope & 252 slots & 3 slice types per snapshot \\
License & CC-BY 4.0 & CC-BY 4.0 \\
\bottomrule
\end{tabular}
\end{table}

We evaluate on two complementary public datasets
(Table~\ref{tab:datasets}).

\subsubsection{Live RAN performance counters (Zenodo 17815388)}
Lehoczky et al.~\cite{lehoczky2026ran} release performance-management
counters from an operational multi-technology network: 302{,}052
sector-level records at 15-minute granularity from 75 sectors across
5G, 4G, and 2G basebands. We aggregate sector data volumes
(DL$+$UL) into a single network demand series, restrict to the
longest gap-free segment (1{,}778 slots, 2023-10-06 to 2023-10-25,
zero missing slots after interpolation of short holes), and set
capacity $C = 143{,}954.5$ (1.3$\times$ the observed peak, the
operator-typical headroom rule). Demand statistics on the segment:
mean 51{,}394, median 59{,}481, maximum 110{,}734 --- right-skewed
and strongly diurnal. The chronological 70/15/15 split yields
1{,}175 training, 251 validation, and 252 test windows.

\subsubsection{CTTC B5G network slicing (Zenodo 10610616)}
Farreras et al.~\cite{farreras2024cttc} publish steady-state
slicing-simulation snapshots, each containing per-slice reservations
(the operator's $\delta$-share of link capacity), offered traffic,
and observed QoS (packet-drop ratio, average delay) for eMBB, mMTC,
and URLLC. We parse the first 250 snapshots in the dataset's native
order; the loader's corrupt-sample guard retains 229, of which 70\%
(by snapshot id) form the training pool for empirical quantiles and
the remainder the evaluation stream. This dataset exercises
Stage~2's policy logic on slice-level, non-time-series data where
reservations and outcomes are directly observable.

\subsection{Baselines}

All baselines commit reservations under the same information
constraints as UCRA and are scored on the same test window and
capacity.

\begin{itemize}
\item \textbf{Static peak} --- reserve the training-window maximum
demand, constant across the test window. The canonical conservative
deployment rule; its failure mode is structural waste.
\item \textbf{Median forecast} --- reserve the Stage-1 median
quantile per slot, i.e., a perfect point forecast with no uncertainty
awareness. Isolates the value of the uncertainty machinery. Scored
against the \emph{frozen} forecast (see
Section~\ref{sec:leakage-audit}).
\item \textbf{Train-only $\Phi$} --- reserve
$q^{\text{train}}_{0.9} + \kappa\,(q^{\text{train}}_{0.99} -
q^{\text{train}}_{0.50})$, a static reservation that applies UCRA's
own kappa-buffer rule but with all statistics frozen from the
training window: no online forecast, no risk feedback, no smoothing.
Isolates what the \emph{adaptive} Stages 1--4 contribute over the
$\Phi$ formula alone. On CTTC, the analogous policy uses the
empirical train-half quantiles per slice type.
\item \textbf{Oracle quantile} --- a cheating upper bound that
computes the rolling 90\% quantile of the \emph{realized} test
demand itself (48-slot window). It cannot be deployed; it answers
``how good is a quantile policy with near-perfect knowledge?'' and
calibrates expectations for what forecast quality can ever deliver.
\end{itemize}

On CTTC we additionally score the \textbf{operator} baseline --- the
dataset's own $\delta$-sized reservations as shipped --- which is the
most policy-relevant comparison of all.

\subsection{Metrics and Protocol}

Reservation metrics follow Section~\ref{sec:system}: violation rate
$V$, utilization $U$, waste $W$, tracking error $E$, computed per
slot (or per snapshot for CTTC) on the test scope. Stage-1
calibration is reported separately as \emph{coverage}: the fraction
of test slots whose realized demand falls inside the forecast
q05--q99 band. Coverage is a property of the forecaster, not of the
reservation policy, and we keep the two layers distinct in all
analysis. Statistical rigor: the seed study
(Section~\ref{sec:results-seeds}) repeats training and evaluation
under seeds $\{41, 42, 43\}$ --- fresh model initialization, fresh
stochastic training, identical evaluation --- and reports mean $\pm$
sample standard deviation. All Stages 2--4 computations are
deterministic given the forecast, so the seed spread measures
exactly the pipeline's dependence on the learned component.

\subsection{Leakage Audit and Causality Regression Tests}
\label{sec:leakage-audit}

A dedicated audit examines every path through which information could
flow from the evaluation future into model or policy decisions. Three
defects were found and fixed during preparation; we document them
because they are instructive and because our final numbers depend on
the fixes.

\emph{L1 --- replay label leakage.} The original Stage-5 harness
offered each test window to the replay buffer immediately at loop
time $t$, together with its full $H$-step target covering slots
$t..t{+}H{-}1$ --- of which $H{-}1 = 3$ (45 minutes) were still in
the future. Any subsequent fine-tune therefore trained on labels no
live system could possess. \textbf{Fix:} the \texttt{DelayedReplay}
staging queue of Section~\ref{sec:stage5} admits a window only at
$t{+}H{-}1$ (Proposition~\ref{prop:causal}).

\emph{L2 --- global forecast refresh.} After an update at time $t$,
the original harness recomputed forecasts for the \emph{entire} test
window. Reservations for slots $\le t$ were already committed, but
the rewrite silently changed (i) the forecast table from which the
median-forecast baseline was computed and (ii) the quantile-band
figure of what was served. \textbf{Fix:} partial refresh ---
only slots $> t$ receive the evolved model's forecasts; served slots
keep their served forecast, and baselines always read the frozen
pre-update forecast.

\emph{L3 --- outcome-informed risk (CTTC).} The original CTTC
evaluation computed its per-sample risk term from the offered load of
the very sample being served --- the policy knew what it was
protecting before committing protection. \textbf{Fix:} a sequential
serving loop in which $\rho$ is computed only from previously served
samples (the same \texttt{UncertaintyState} contract as the RAN
loop), and the current sample is observed only after its reservation
is fixed.

Regression tests in \texttt{tests/test\_causality.py} enforce: (i) a
window offered at $t$ is absent from the buffer at $t{+}H{-}2$ and
present at $t{+}H{-}1$; (ii) any fine-tune batch contains only labels
$\le$ the update time; (iii) a refresh at $t$ leaves entries
$\le t$ bit-identical. The tests fail against the leaky variants and
pass against the final pipeline. In addition, all three defects
existed in the pre-submission pipeline and the corrected headline
numbers \emph{changed}: the drift-demo gap narrowed modestly (the
leak had inflated adaptation gains), and the CTTC violations rose
slightly --- the honest direction in both cases.

\subsection{Implementation and Runtime}

The pipeline is pure Python (PyTorch 2.x CPU, NumPy, pandas) with
YAML-driven configuration; no GPU is required at any stage. Stage-1
training completes in under three minutes per seed on two CPU cores;
the 252-slot closed loop runs in seconds; the drift demo (two passes
plus up to three fine-tunes) in under a minute; the CTTC evaluation
parses 250 snapshots in $\approx$17 minutes single-threaded
($\approx$7 minutes with the two-worker archive-parallel driver).
The three-seed study (train + loop + sweep + demo per seed, plus
aggregation) completes in a single ten-minute invocation. Every
figure and table in this paper is produced by committed scripts from
the JSON/CSV artifacts the pipeline writes.

\emph{Released artifacts.} The repository ships: the \texttt{ucra}
package (loaders, windowing, model, stages, metrics, baselines,
plots); the six experiment drivers used for this paper
(\texttt{train}, \texttt{run\_ucra}, \texttt{demo\_evolution},
\texttt{run\_cttc\_eval}, \texttt{run\_cttc\_parallel},
\texttt{run\_seeds}) plus the data audit and figure scripts; the two
test suites (\texttt{test\_smoke.py}, \texttt{test\_causality.py});
and this paper's \LaTeX{} source with the verified result JSONs, so
that tables can be cross-checked against the exact numbers that
produced them.
